\chapter{3. Propagation speed and kernel estimates}
\source{cheeger-gromov-taylor:page-0020,cheeger-gromov-taylor:page-0021,cheeger-gromov-taylor:page-0022,cheeger-gromov-taylor:page-0023,cheeger-gromov-taylor:page-0024,cheeger-gromov-taylor:page-0025,cheeger-gromov-taylor:page-0026,cheeger-gromov-taylor:page-0027}

The estimate (2.13) on $k_f(x_1,x_2)$ for $f\in \mathscr{F}_2^{\varphi,A}$ holds with $A=d/2$, and we have

\begin{equation}\tag{3.4}
|k_f(x_1,x_2)|\le c_n^2\left(\frac{2}{a}\right)^n \psi(r)^2,
\end{equation}

if

\begin{equation}\tag{3.5}
\rho(x_1,x_2)=a+r.
\end{equation}

If we do not assume a lower bound on the Ricci curvature, methods of estimating the gradient of $k_f(x_1,x_2)$ do not apply. We note that the assumption of $C^0$-bounded geometry implies a H\"older bound on $k_f(x_1,x_2)$ which improves the pointwise bound (3.4) for $f\in \mathscr{F}_2^{\varphi,A}$.

Indeed, according to Corollary 1.2 if $u$ is supported on $\Ba(x_2)$, then

\begin{equation}\tag{3.6}
\bigl\|\Delta^k f(\sqrt{-\Delta}) \Delta^l u\bigr\|_{M\setminus B_{r+a}(x_2)}
\le \frac{1}{\pi} C(C(2k+2l))^{2k+2l}\psi(r)\|u\|.
\end{equation}

Thus if $\rho(x_1,x_2)=r+2a$, then

\begin{equation}\tag{3.7}
\bigl\|\Delta^k f(\sqrt{-\Delta}) \Delta^l u\bigr\|_{B_a(x_1)}
\le \frac{C}{\pi}(C(2k+2l))^{2k+2l}\psi(r)\|u\|.
\end{equation}

We now argue similarly to the proof of Theorem 2.1. Given $u$ supported in $B_a(x_2)$, if $|y|\le B=\tfrac12 c^{-1}$, the power series

\begin{equation}\tag{3.8}
v(y,x)=\sum_{k=0}^{\infty}\frac{y^{2k}}{(2k)!}\,\Delta^k f(\sqrt{-\Delta})u
\end{equation}

converges to an element of $L^2([-B,B]\times B_a(x_1))$, and we have

\begin{equation}\tag{3.9}
\left(\frac{\partial^2}{\partial y^2}+\Delta\right)v=0,
\end{equation}

\begin{equation}\tag{3.10}
\|v\|_{[-B,B]\times B_a(x_1)}\le C_1\psi(r)\|u\|_{L^2}.
\end{equation}

Now in normal coordinates on $B_a(x_1)$ we have

\begin{equation}\tag{3.11}
\Delta w(x)=\sum_{j,k} g^{-1/2}\frac{\partial}{\partial x_j}
\left(g^{1/2}g^{ij}\frac{\partial w}{\partial x_k}\right),
\end{equation}

and our hypothesis that $M$ has $C^0$-bounded geometry precisely gives uniform ellipticity of (3.11), and hence of (3.9), with $C^0$ bounds on the coefficients. Thus the di Georgi-Nash-Moser theorems on such elliptic equations in divergence form (see Gilbarg and Trudinger [20]) imply the pointwise estimate

\begin{equation}\tag{3.12}
|v(x_1)|\le C_2\psi(r)\|u\|
\end{equation}

\[
\text{for }u\text{ supported in }\Ba(x_2).\ \text{Now }(3.12)\text{ is equivalent to the assertion that}
\]
\[
f(\sqrt{-\Delta})\delta_{x_1}\text{ is }L^2\text{ on }B_a(x_2)\text{ and}
\]
\begin{equation}\tag{3.13}
 \bigl\|f(\sqrt{-\Delta})\delta_{x_1}\bigr\|_{B_a(x_2)}\le c_2\psi(r).
\end{equation}

\[
\text{Similarly, we obtain}
\]
\begin{equation}\tag{3.14}
 \bigl\|\Delta^l f(\sqrt{-\Delta})\delta_{x_1}\bigr\|_{B_a(x_2)}\le c_2(2Cl)^{2l}\psi(r).
\end{equation}

\[
\text{So we can set, for }|y|\le B,
\]
\begin{equation}\tag{3.15}
 v_{x_1}(y,x)=\sum_{l=0}^{\infty}\frac{y^{2l}}{(2l)!}\Delta^l f(\sqrt{-\Delta})\delta_{x_1}
\end{equation}
\[
\text{to obtain}
\]
\begin{equation}\tag{3.16}
 \left(\frac{\partial^2}{\partial y^2}+\Delta\right)v_{x_1}=0,
\end{equation}
\begin{equation}\tag{3.17}
 \bigl\|v_{x_1}\bigr\|_{[-B,B]\times B_a(x_2)}\le c_3\psi(r).
\end{equation}

\[
\text{Again the De Giorgi-Nash-Moser results yield a bound on }v_{x_1}(x_2),\text{ i.e.,}
\]
\begin{equation}\tag{3.18}
 k_f(x_1,x_2)\le c_4\psi(r).
\end{equation}

\[
\text{But in fact the theory yields a Hölder bound}
\]
\begin{equation}\tag{3.19}
 \bigl\|k_f(x_1,\cdot)\bigr\|_{C^\alpha(B_{a/2}(x_2))}\le C_5\psi(r).
\end{equation}

\[
\text{for a certain }\alpha>0\text{ which depends only on the }C^0\text{-bound and ellipticity}
\]
\[
\text{constant in (3.11). To summarize, we have proved the following.}
\]

\begin{theorem}[Kernel estimates for $C^0$-bounded geometry]
\label{thm:c0-bounded-geometry-kernel-estimates}
\dcref{Theorem 3.1}
\group{section-3}
\source{cheeger-gromov-taylor:page-0020,cheeger-gromov-taylor:page-0021}
\uses{cor:derivative-tail-estimate}
\textit{If $M$ has $C^0$-bounded geometry and $f\in \mathcal{F}_2^{\varphi,A}$, then the kernel $k_f(x_1,x_2)$ satisfies the estimate (3.19) and also}
\begin{equation}\tag{3.20}
 \bigl\|k_f(\cdot,x_2)\bigr\|_{C^\alpha(B_{a/2}(x_1))}\le C_5\psi(r).
\end{equation}
\end{theorem}

\[
\text{The fact that (3.20) holds follows from the symmetry of }k_f(x_1,x_2).
\]
\[
\text{In particular, recall from the appendix to \S2 that if }f(\lambda)\text{ is the restriction to}
\]
\[
\lambda\in\mathbf{R}\text{ of a function }f(z)\text{ holomorphic on a region}
\]
\begin{equation}\tag{3.21}
 \Omega=\{z\in \mathbf{C}:\ |\operatorname{Im}z|\le W+\epsilon\}\cup\{z\in \mathbf{C}:\ |\operatorname{Im}z|\le B\,|\operatorname{Re}z|\},
\end{equation}
\[
\text{then }f\in \mathcal{F}_2^{\varphi,A}\text{ with }\varphi(s)=Ce^{-Ws},\text{ provided}
\]
\begin{equation}\tag{3.22}
 |f(z)|\le C(1+|z|)^m,\qquad z\in\Omega.
\end{equation}

\[
\text{This gives}
\]

\begin{corollary}[Exponential kernel decay under holomorphy]
\label{cor:c0-bounded-geometry-exponential-decay}
\dcref{Corollary 3.2}
\group{section-3}
\source{cheeger-gromov-taylor:page-0021,cheeger-gromov-taylor:page-0022}
\uses{thm:c0-bounded-geometry-kernel-estimates}
\textit{If $M$ has $C^0$-bounded geometry, and $f(z)$ is holomorphic on (3.21), satisfying (3.22), then}
\begin{equation}\tag{3.23}
 |k_f(x_1,x_2)|\le Ce^{-Wr},
\end{equation}
\end{corollary}

\[
\text{provided}
\]
\begin{equation}\tag{3.24}
\rho(x_1,x_2)=r+2a.
\end{equation}

We will not pursue the implication of $C^k$-bounded geometry for finite positive $k$, but will devote the rest of this section to obtaining refined global results for manifolds with $C^\infty$ or $C^\omega$ bounded geometry.

We now consider the case when $M$ has $C^\infty$-bounded geometry. Examples include all homogeneous spaces, i.e., manifolds with a transitive group of isometries, as well as more general sorts of manifolds such as leaves of $C^\infty$ foliations of compact manifolds. Since such $M$ has bounded curvature, the estimates of \S1 apply to the estimate of the kernel $k_f(x_1,x_2)$ for $\rho(x_1,x_2)\ge a$. We complement this with a precise analysis of the behavior of the kernel of $f(\sqrt{-\Delta})$ for $\rho(x_1,x_2)\le a$. Assume $f(\lambda)\in S_{\rho,0}^m(\mathbf{R})$, i.e.,

\begin{equation}\tag{3.25}
|f^{(j)}(\lambda)|\le C_j(1+|\lambda|)^{m-j\rho}.
\end{equation}

We use here and below the notation for symbol classes and pseudodifferential operators of H\"ormander; see [32, Chapter II]. From the representation

\begin{equation}\tag{3.26}
f(\sqrt{-\Delta})u=\frac{1}{2\pi}\int_{-\infty}^{\infty}\hat f(s)\cos s\sqrt{-\Delta}\,u\,ds,
\end{equation}

use a partition of unity $1=\phi_1(s)+\phi_2(s)$, $\phi_1(s)\in C_0^\infty(-a,a)$, even, $\phi_1(s)=1$ for $|s|\le a/2$, to get

\begin{align}
f(\sqrt{-\Delta})u
&=\frac{1}{2\pi}\int_{-a}^{a}\phi_1(s)\hat f(s)\cos s\sqrt{-\Delta}\,u\,ds \notag\\
&\quad +\frac{1}{2\pi}\int_{-\infty}^{\infty}\phi_2(s)\hat f(s)\cos s\sqrt{-\Delta}\,u\,ds \notag\\
&=F_1u+F_2u.
\tag{3.27}
\end{align}

Note that

\begin{equation}\tag{3.28}
F_2u=(1-\Delta)^{-k}\int_{-\infty}^{\infty}
\left(1-\frac{d^2}{ds^2}\right)^k(\phi_2(s)\hat f(s))\cos s\sqrt{-\Delta}\,u\,ds
=(1-\Delta)^{-k}G_ku,
\end{equation}

with $\|G_ku\|\le C_k\|u\|$. Thus the standard elliptic estimates and the Sobolev embedding theorem give

\begin{equation}\tag{3.29}
\|F_2u\|_{C^l(B_a(x_1))}\le C_l\|u\|_{L^2}.
\end{equation}

Furthermore, if $u$ is supported in $B_a(x_1)$,

\begin{equation}\tag{3.30}
\|F_2u\|_{C^l(B_a(x_1))}\le C_{lk}\|u\|_{H^{-k}(B_a(x_1))},
\end{equation}

where the norm on the right is a Sobolev norm.

Now consider
\begin{equation}\tag{3.31}
F_1u=\frac{1}{2\pi}\int_{-a}^{a}\varphi_1(s)\hat f(s)\cos s\sqrt{-\Delta}\,u\,ds,
\end{equation}
with $u$ supported in $B_a(x_1)$. We will analyze (3.31) by replacing the operator $\cos s\sqrt{-\Delta}$, for $|s|\le a$, by its geometrical optics approximation. This approach to the study of $f(\sqrt{-\Delta})$ on compact manifolds is emphasized in Chapter XII of [32], and in [33], and also plays a role in the work of [8] on the wave equation on a cone.

By shrinking $a$ if necessary, we can assume the parametrix for $\cos s\sqrt{-\Delta}$ is of the following form, in a normal coordinate system centered at $x_1$:
\begin{equation}\tag{3.32}
\cos s\sqrt{-\Delta}\,u(x)=\sum_{j=1}^{2}\int a_j(s,x,\xi)e^{i\varphi_j(s,x,\xi)}\hat u(\xi)\,d\xi.
\end{equation}
Here $\varphi_j$ solves the eikonal equation
\begin{equation}\tag{3.33}
\frac{\partial}{\partial s}\varphi_j = (-1)^j\lambda_1^j(x,d\varphi_j)=(-1)^j(d\varphi_j,d\varphi_j)^{\frac12},
\end{equation}
\begin{equation}\tag{3.34}
\varphi_j(0,x,\xi)=x\cdot \xi
\end{equation}
for $|s|\le a$, $x\in B_{2a}(x_1)$; $a_j(s,x,\xi)$ is determined by the usual transport equations of geometrical optics (see for example, [32, Chapter VIII]), with $a_j(0,x,\xi)=\frac12$. Now if $u$ is supported in $B_a(x_1)$, then, for $|s|\le a$, $\cos s\sqrt{-\Delta}\,u$ is supported in $B_{2a}(x_1)$. If we set $\hat g(s)=\varphi_2(s)\hat f(s)$, then $g(\lambda)\in S^{m}_{\rho,0}$ differs from $f(\lambda)$ by an element of the Schwartz space $\mathcal S(\mathbb R)$ of rapidly decreasing functions. Thus we have
\begin{equation*}
F_1u=\sum_{j=1}^{2}\int_{-a}^{a}\hat g(s)a_j(s,x,\xi)e^{i\varphi_j(s,x,\xi)}\hat u(\xi)\,d\xi\,ds+S_1u
\end{equation*}
\begin{equation}\tag{3.35}
=\sum_j\int g(D_s)\!\left(a_j(s,x,\xi)e^{i\varphi_j(s,x,\xi)}\right)\Big|_{s=0}\hat u(\xi)\,d\xi+S_1u
\end{equation}
\begin{equation*}
=\int b(x,\xi)e^{ix\cdot \xi}\hat u(\xi)+S_1u,
\end{equation*}
where $S_1:\mathcal E'(B_a(x_1))\to C^\infty(B_{3a}(x_1))$, and $b(x,\xi)$ is given by the fundamental asymptotic expansion lemma for pseudodifferential operators. Thus
\begin{equation}\tag{3.36}
b(x,\xi)=\sum_{j=1}^{2}b_j(0,x,\xi),
\end{equation}
where
\begin{equation}\tag{3.37}
b_j(s,x,\xi)=e^{-i\varphi_j(s,x,\xi)}g(D_s)\!\left(a_j(s,x,\xi)e^{i\varphi_j(s,x,\xi)}\right).
\end{equation}

belongs to $S^m_{\rho,1-\rho}$ if $\rho \ge \frac12$ with
\begin{equation}\tag{3.38}
b_j(s,x,\xi)\sim a_j(s,x,\xi)g\!\left(\frac{\partial}{\partial s}\varphi_j\right)+\cdots.
\end{equation}
Hence using (3.33) and (3.34) we obtain
\begin{equation}\tag{3.39}
b(x,\xi)\sim g(\lambda_1(x,\xi))+\cdots.
\end{equation}
The remainder terms in (3.38) and (3.39) sum to elements of $S^{m-(2\rho-1)}_{\rho,1-\rho}$ for $\rho \ge \frac12$ and we have a complete asymptotic expansion for $\rho>\frac12$.

Denoting by $\OPS^m_{\rho,\delta}$ the set of pseudodifferential operators with symbols in the class $S^m_{\rho,\delta}$, we have the following conclusion.

\begin{theorem}[Pseudodifferential representation of $f(\sqrt{-\Delta})$]
\label{thm:pdo-representation-of-functional-calculus}
\dcref{Theorem 3.3}
\group{section-3}
\source{cheeger-gromov-taylor:page-0023,cheeger-gromov-taylor:page-0024}
\textit{Assume $M$ has $C^\infty$-bounded geometry. Let $f(\lambda)\in S^m_\rho(\R)$, $\rho \ge \frac12$ be even. If $u$ is supported in $\Ba(p)$ in normal coordinates centered at $p$, then we have}
\begin{equation}\tag{3.40}
f(\sqrt{-\Delta})u=B_p(x,D)u;\ x\in \Ba(0),
\end{equation}
\textit{where $B_p(x,D)\in \OPS^m_{\rho,1-\rho}$, and in fact}
\begin{equation}\tag{3.41}
\{B_p(x,D): p\in M\}\text{ is bounded in }\OPS^m_{\rho,1-\rho}.
\end{equation}
\textit{If $\rho>\frac12$, the symbol of $B_p(x,D)$ has a complete asymptotic expansion of the form (3.39), valid uniformly for $p\in M$.}
\end{theorem}

We remark that it would be quite natural to use a Hadamard type parametrix written invariantly in terms of the Riemannian distance instead of the Fourier integral representation (3.32), but we will not pursue this approach.

We can use Theorem 3.3 together with the estimates on the kernel $k_f(x_1,x_2)$ for $\rho(x_1,x_2)\ge a$ which follow as a special case of the results of \S1, to obtain $L^p$-norm estimates on $f(\sqrt{-\Delta})$ for the following class of function $f$.

\begin{definition}[Holomorphic symbol class on a strip]
\label{def:holomorphic-strip-symbol-class}
\dcref{Definition 3.1}
\group{section-3}
\source{cheeger-gromov-taylor:page-0024}
We say $f\in S_w^m$ if $f(z)$ is holomorphic on the strip
\begin{equation}\tag{3.42}
\Omega_w=\{z\in \C : |\Impart z|<W\},
\end{equation}
and satisfies the $S^m_{1,0}$ symbol estimates on $\Omega_w$:
\begin{equation}\tag{3.43}
|f^{(j)}(z)|\le C_j(1+|z|)^{m-j},\qquad z\in \Omega_w.
\end{equation}
\end{definition}

We have, according to (3.27),
\begin{equation}\tag{3.44}
f(\sqrt{-\Delta})=\Phi_1+\Phi_2,
\end{equation}
where the kernel of $\Phi_1$ is supported in the neighborhood $\{(x_1,x_2): \rho(x_1,x_2)\le 2a\}$ of the diagonal in $M\times M$ and is given in local normal coordinates as a bounded family of pseudodifferential operators. The kernel $\Phi_2(x_1,x_2)$ is smooth on $M\times M$ and satisfies the estimate
\begin{equation}\tag{3.45}
|\Phi_2(x_1,x_2)|\le Ce^{-Wr},\qquad r=\rho(x_1,x_2).
\end{equation}

From this we can deduce $L^{p}$ continuity results for $f(\sqrt{-\Delta})$, as follows. First we show $\Phi_1: L^{p}(M)\to L^{p}(M)$, $1<p<\infty$. Let $u\in L^{p}(M)$. We can write

\begin{equation}\tag{3.46}
u=\sum_{j=0}^{\infty}u_{j}
\end{equation}

where each $u_{j}$ is supported in a ball $B_{a}(p_{j})$ of radius $a$ in $M$, and, because we are assuming $M$ has bounded geometry, we can suppose there is a constant $K$ such that, for any $j$, $B_{2a}(p_{j})$ intersects no more than $K$ other balls $B_{2a}(p_{k})$, $p_{k}\neq p_{j}$. Thus we can arrange that

\begin{equation}\tag{3.47}
C^{-1}\sum_{j}\|u_{j}\|_{L^{p}}\leq \|u\|_{L^{p}}\leq C\sum_{j}\|u_{j}\|_{L^{p}}.
\end{equation}

Now $\Phi_{1}u=\sum_{j}\Phi_{1}u_{j}$, and operators $b(x,D)$ with symbols in $S^{0}_{1,0}$ are continuous on $L^{p}$, $1<p<\infty$, and bounded families of such operators have uniformly bounded operator norm on $L^{p}$, $1<p<\infty$; see [32, Chapter XII]. Thus from (3.40) and (3.41) it follows that

\begin{equation}\tag{3.48}
\|\Phi_{1}u_{j}\|_{L^{p}}\leq C_{1}\|u_{j}\|_{L^{p}},
\end{equation}

provided $f\in S^{0}_{1,0}(\mathbf{R})$. Furthermore, $\Phi_{1}u_{j}$ is supported in $B_{2a}(p_{j})$. It follows that

\begin{equation}\tag{3.49}
\|\Phi_{1}u_{j}\|_{L^{p}}\leq C_{2}(p)\|u\|_{L^{p}}, \qquad \text{if } f\in S^{0}_{1,0}(\mathbf{R}).
\end{equation}

As for the $L^{p}$ continuity of $\Phi_{2}$, we use the estimate (3.45) together with the elementary fact that the $L^{p}$ operator norm of $\Phi_{2}$ is bounded by

\begin{equation}\tag{3.50}
\sup_{q}\int_{M}|\Phi_{2}(p,q)|\,d\vol(q)+\sup_{q}\int_{M}|\Phi_{2}(p,q)|\,d\vol\,p.
\end{equation}

See [32, Chapter XIII]. Now if $M$ has $C^{\infty}$-bounded geometry, or more generally if $M$ has Ricci curvature bounded below, we have

\begin{equation}\tag{3.51}
\vol\{x\in M:\rho(x,p)\leq r\}\leq ce^{Kr}
\end{equation}

for some $K$ independent of $p$. Consequently, if (3.45) holds with $W>K$, we have a uniform bound on (3.50) and hence an $L^{p}$-operator bound on $\Phi_{2}$, $1\leq p\leq \infty$. Thus we have proved the following result.

\begin{theorem}[\texorpdfstring{$L^p$}{Lp} continuity for strip-holomorphic symbols]
\label{thm:lp-continuity-strip-holomorphic-symbols}
\dcref{Theorem 3.4}
\group{section-3}
\source{cheeger-gromov-taylor:page-0025}
\textit{Assume $M$ has $C^{\infty}$ bounded geometry. If}
\begin{equation}\tag{3.52}
f\in \Swo, \qquad W>K,
\end{equation}
\textit{where (3.51) holds for $K$, then}
\begin{equation}\tag{3.53}
f(\sqrt{-\Delta}): L^{p}(M)\to L^{p}(M), \qquad 1<p<\infty.
\end{equation}
\end{theorem}

At this point we make the following remark. Sometimes one knows that the spectrum of $-\Delta$ on $M$ is contained in $[\alpha_{0},\infty)$ for some $\alpha_{0}>0$. In such a case, the same considerations as above, or indeed elsewhere in this paper, apply to

\begin{equation}\tag{3.54}
 f(\sqrt{-\Delta}-\alpha_0): L^p(M) \to L^p(M), \qquad 1 < p < \infty,
\end{equation}

for $f\in \mathcal S_w^0$. Some results a bit weaker than Theorem 3.4 were announced in [33].

We pass now to a still more rigid class of manifolds, those with $C^\omega$-bounded geometry. This class still includes homogeneous spaces. In this case, if $f\in \mathcal F_2^{\varphi,\mathcal A}$, the construction (3.8) yields a harmonic function on $(-B,B)\times B_a(x_1)$, and also (3.15) yields a harmonic function on $(-B,B)\times B_a(x_2)$, i.e., these functions are annihilated by the operator $\partial^2/\partial y^2 + \Delta$, and in normal coordinates we have $\Delta$ of the form (3.11). This time the coefficients of $\partial^2/\partial y^2 + \Delta$ form a bounded family of analytic functions on $B_a(0)$ with a uniform ellipticity constant, so the analyticity estimates for solutions to $(\partial^2/\partial y^2 + \Delta)w = 0$ imply the estimates on the kernel $k_f(x_1,x_2)$:

\begin{equation}\tag{3.55}
\sup_{x\in B_a(x_1)} |D_x^\alpha k_f(x_1,x_2)| \le C(C|\alpha|)^{|\alpha|}\psi(r),
\end{equation}

if $\rho(x_1,x_2)=r+2a$, where one uses normal coordinates centered at $x_1$ to express $D_x^\alpha$. In particular, if $f(z)$ is holomorphic in the region

\begin{equation}\tag{3.56}
\Omega = \{z\in C: |\Im z|\le W+\varepsilon\}\cup\{z\in C: |\Im z|\le B\, |\Re z|\}
\end{equation}

and satisfies

\begin{equation}\tag{3.57}
|f(z)|\le C(1+|z|)^m, \qquad z\in \Omega,
\end{equation}

then, as discussed in the appendix to \S2, the estimate (3.55) becomes

\begin{equation}\tag{3.58}
\sup_{x\in B_a(x_1)} |D_x^\alpha k_f(x_1,x_2)| \le C(C|\alpha|)^{|\alpha|} e^{-rW}.
\end{equation}

Since the function $f_t(\lambda)=e^{-t\lambda^2}$ belongs to $\mathcal F_\omega^{\varphi,\mathcal A}$ with $\varphi(r)=e^{-(1-\varepsilon)r^2/4t}$, we also obtain from (3.55) estimates on the heat kernel proved by G\aa rding [19], for the case where $M$ is a Lie group with left invariant metric.

Finally, we note a dampening of the kernel $k_f(x_1,x_2)$ when $M$ is a rank 1 symmetric space, i.e., when the isotropy group of $M$ fixing a point $x$ acts transitively on the unit sphere in $T_xM$. We suppose $M$ is of noncompact type with negative curvature. In such a case, the estimate (3.51) can be refined to the following sharper result on the volume of a spherical shell

\begin{equation}\tag{3.59}
A_r(q) = \{x\in M: a+r\le \rho(q,x)\le 3a+r\}.
\end{equation}

Namely,

\begin{equation}\tag{3.60}
\operatorname{vol} A_r(q) \sim C(a)e^{Kr},
\end{equation}

which implies
\begin{equation}\tag{3.61}
\operatorname{vol}\{x \in M : \rho(x,q) \leq r + 2a\} \sim Ce^{Kr}.
\end{equation}
Now $f(\sqrt{-\Delta})$ is invariant under rotations of $M$ about $q$. This fact together with (3.60) and the usual elliptic estimates yields
\begin{equation}\tag{3.62}
\lvert k_f(x,q)\rvert \leq Ce^{-Kr/2} \sup_{0<k\leq [n/2] + 1}
\int_r^\infty \lvert \hat{f}^{(2k)}(s)\rvert\, ds,
\end{equation}
if
\begin{equation}\tag{3.63}
\rho(x,q)=r+2a.
\end{equation}
In particular, by (3.60) the function $u_q(x)=k_f(x,q)$ is integrable on $\{x \in M :
\rho(x,q)\geq 1\}$ with $L^1$ norm independent of $q$, provided
\begin{equation}\tag{3.64}
\int_r^\infty \lvert \hat{f}^{(2k)}(s)\rvert\, ds \leq Cr^{-1-\epsilon}e^{-Kr/2},\qquad
0\leq k\leq \left[\frac{n}{2}\right]+1.
\end{equation}
It follows that, for a rank-one symmetric space $M$, if (3.66) holds,
\begin{equation}\tag{3.65}
f(\sqrt{-\Delta}) : L^p(M)\to L^p(M),\qquad 1<p<\infty.
\end{equation}
More generally, if the spectrum of $-\Delta$ is contained in $[\alpha_0,\infty)$, and
\begin{equation}\tag{3.66}
f\in \mathcal{S}_w^0,\qquad W>\frac{1}{2}K,
\end{equation}
then
\begin{equation}\tag{3.67}
f\!\left(\sqrt{-\Delta-\alpha_0}\right): L^p(M)\to L^p(M),\qquad 1<p<\infty.
\end{equation}
This result was proved by Stanton and Tomas [31] by a different method, following the work of Clerc and Stein [12] on complex symmetric spaces. As these authors point out, one can make use of the Kunze-Stein phenomenon which implies that the operator $\Phi_2$ of (3.44) is continuous on $L^p(M)$, $1<p<$
$\infty$, provided that, for some $b>1$,
\begin{equation}\tag{3.68}
\Phi_2(x_1,\cdot)\in L^q(M),\qquad q\in(1,b).
\end{equation}
which is slightly weaker than the requirement
\begin{equation}\tag{3.69}
\Phi_2(x_1,\cdot)\in L^1(M),
\end{equation}
we sought above. Indeed, this allows us to replace $W>\frac{1}{2}K$ by
\begin{equation}\tag{3.70}
W=\frac{1}{2}K
\end{equation}
in (3.66).
